\section{Related Work}\label{sec:related}

A word first on how we use the literature.  The active-inference citations in this paper are implementation warrants rather than background gloss: the quantities that drive the harness are drawn from a small working canon --- Friston et al.\ for the canonical decomposition of expected free energy into risk and ambiguity and for the treatment of habit \cite{FRISTON2016862}, the free-energy principle paper for the underlying formulation \cite{friston2019freeenergyprincipleparticular}, and Smith et al.\ for the discrete-state account of concept learning \cite{smith2020structure} --- and each Glossary entry cites the source that fixes the meaning of its quantity while footnoting the code that implements it.  Victorita Oana Neacsu's PhD thesis on ``Structure Learning and learning structure with Active Inference'' \cite{neacsu2024structure} plays a particular role: it describes the strategy that we implement to update our grasp of the pattern landscape after each round --- Bayesian Model Reduction over Dirichlet counts, with an explicit evidence threshold (see the Glossary).  Where the implementation departs from the formalism, the departure is named in place rather than smoothed over; the Glossary's notational warning about engineering scores versus free energies is one instance of this discipline.

The intersection of active inference and language-model agents is young and, so far, dominated by position papers.  Pezzulo, Friston, and colleagues argue that generative AI as currently deployed is ``passive'' --- lacking the action--perception loop that grounds meaning \cite{pezzulo2024generating} --- and Kulveit et al.\ read LLMs as atypical active-inference agents precisely because the tight act--perceive feedback loop is missing \cite{kulveit2023predictive}.  Implemented systems in which expected free energy actually selects actions for an LLM-based agent remain rare: Prakki's multi-LLM system, where a discrete generative model selects prompt and search strategies \cite{prakki2024active}, and Beckenbauer et al.'s Orchestrator, which uses active inference to monitor multi-agent dynamics in long-horizon tasks \cite{beckenbauer2025orchestrator}, are the closest we have found.  On the industrial side, a recently granted patent --- with Friston among the inventors --- covers the compilation of natural-language specifications into active-inference agents over a knowledge graph \cite{verses2025patent}.  None of this work addresses software engineering.

Agentic coding systems, meanwhile, generally close their loops informally.  Reflexion, SWE-agent, and OpenHands iterate propose--test--reflect cycles without a decision-theoretic objective \cite{shinn2023reflexion,yang2024sweagent,wang2025openhands}; a self-improving coding agent can hill-climb its own codebase against a benchmark utility \cite{robeyns2025sica}; SWE-Search adds Monte-Carlo tree search with a learned value function \cite{antoniades2024swesearch}.  A recent position paper diagnoses exactly this gap --- engineering practice has tool loops and memory, but ``lacks principled grounding'' --- and calls for a cybernetics of foundation agents \cite{wang2026cybernetics}.  The nearest neighbour to the present work is Papamarkou et al.'s Bayesian controller for coding agents, which maintains a belief over candidate-code correctness and decides whether to gather evidence, refine, verify, or stop by value of information in a partially observed decision process \cite{papamarkou2026bayesian}.  Value of information is, in effect, the epistemic term of expected free energy; the differences are that active inference carries preference and information gain in a single objective over a declared generative model, learns precision from realised outcomes, and --- in this system --- spans strategic mission selection through enactment and structure learning, rather than controlling the refine--verify--stop cycle of a single candidate program.

To our knowledge, then, as of mid-2026 no previously published system drives an agentic coding loop with the active-inference formalism proper --- generative model, expected free energy, precision, and structure learning together.  We state this as a landscape observation rather than a priority contest: the space is moving quickly, the surrounding literature closes its loops either informally or with adjacent decision-theoretic machinery, and our response to the nearest of these is engagement rather than avoidance.

